% Generated draft: requires mathematical and source review.
\documentclass[11pt]{article}
\usepackage[margin=25mm]{geometry}
\usepackage{fontspec}
\setmainfont{Latin Modern Roman}
\usepackage{amsmath,amssymb,mathtools,unicode-math}
\setmathfont{Latin Modern Math}
\usepackage{enumitem,needspace,xurl,hyperref}
\setlength{\parindent}{0pt}
\setlength{\parskip}{5pt}
\setlength{\emergencystretch}{3em}
\newcommand{\N}{\mathbb{N}}
\newcommand{\Z}{\mathbb{Z}}
\newcommand{\Q}{\mathbb{Q}}
\newcommand{\R}{\mathbb{R}}
\newcommand{\C}{\mathbb{C}}
\newcommand{\F}{\mathbb{F}}
\newcommand{\A}{\mathbb{A}}
\newcommand{\PP}{\mathbb{P}}
\newcommand{\E}{\mathbb{E}}
\newcommand{\eps}{\varepsilon}
\newcommand{\dd}{\,\mathrm{d}}
\begin{document}
\section*{UnsolvedMath ID: 701361}
\textbf{Source problem code:} catalogue entry
\subsection{Definitions and mathematical statement}
Let \(\operatorname{Sym}(\Z)\) denote the group of bijections from \(\Z\) to itself, and put \(\Z_{\geq 0}=\{0,1,2,\ldots\}\). A permutation \(f\in\operatorname{Sym}(\Z)\) is \emph{residue-class-wise affine} if there exist an integer \(m\geq 1\) and integers \(a_r\neq 0\), \(b_r\), for each \(r\in\{0,\ldots,m-1\}\), such that
\[
 f(r+mk)=a_rk+b_r\qquad\text{for every }k\in\Z.
\]
This expresses the usual meaning of affine on each residue class: equivalently, on that class the formula is \(f(n)=(a_r/m)n+b_r-a_rr/m\), with rational coefficients. Bijectivity is required globally.

Define
\[
 G_+=\bigl\{f\in\operatorname{Sym}(\Z): f\text{ is residue-class-wise affine and }f(\Z_{\geq0})=\Z_{\geq0}\bigr\}.
\]
The setwise condition does not require the nonnegative integers to be fixed individually.

The supplied question, proposed by S. Kohl in 2010 as Kourovka Notebook Problem 17.59, asks whether
\[
 \mathrm{CT}(\Z)=G_+.
\]
The supplied record does not define \(\mathrm{CT}(\Z)\). Thus this equality preserves the stated target, but a complete mathematical formulation requires the intended definition of its left-hand side. The registry labels the problem open; that label is supplied metadata, not an independently verified assessment.
\subsection{Short English statement}
S. Kohl asks whether \(\mathrm{CT}(\Z)\) consists exactly of the permutations of the integers that are affine on each residue class for some common positive modulus and preserve the nonnegative integers setwise. The supplied record leaves \(\mathrm{CT}(\Z)\) undefined, so the intended group must be specified before the question is fully precise.
\subsection{Sources}
\begin{itemize}
\item Supplied local UnsolvedMath registry record, permanent ID \(701361\), exact source problem code \texttt{catalogue entry}, titled \emph{Kourovka Notebook Problem 17.59}. The supplied status is \texttt{open}.
\item Attributed source: \emph{The Kourovka Notebook (Unsolved Problems in Group Theory)}, 17th issue (2010), Problem 17.59; proposed by S. Kohl in 2010. The background also identifies the 21st issue edition (2026), \texttt{arXiv:1401.0300v48}.
\item The supplied collection description attributes the notebook, Issues 1--21 (1965--2026), to editors E. I. Khukhro and V. D. Mazurov. These bibliographic details and the status are supplied attribution and have not been independently verified here.
\end{itemize}
\end{document}
